\section{Introduction}

Student dropout in higher education continues to impose substantial costs on individuals, families, institutions, and society. Higher-education leaders have emphasized the persistence of the problem, arguing that the core challenge is not only financing but also completion and persistence \cite{Crow2023CompletionCrisis,Shapiro2023SomeCollegeNoCredential}.

A large body of work has studied dropout prediction using machine-learning models built from demographic, academic, and engagement signals. However, many operational approaches still emphasize static, single-time risk estimates that primarily address \emph{who} is likely to drop out, offering limited resolution about \emph{when} risk intensifies over the course timeline \cite{Prenkaj2021SurveyMLDropout,AndradeGiron2023DropoutMLReview}. This mismatch often constrains institutions to reactive support or poorly timed interventions relative to the onset of disengagement. In practice, moving from prediction to action requires (i) expressing risk in temporal terms, (ii) translating temporal risk into auditable decision rules and timing, and (iii) comparing alternative support scenarios as explicit intervention regimes.

\noindent\textit{From prediction to policy.}
Static risk scores primarily indicate \emph{who} is at risk, but operational support also depends on \emph{when} risk increases and which decision rule triggers action. Once risk is expressed as a time-indexed hazard trajectory, the next step is to compare explicit intervention rules as regime scenarios (e.g., triggering criteria and short windows of hypothetical hazard reduction). Here, the policy layer is used for \emph{structural} evaluation: it contrasts survival trajectories implied by the model under alternative regimes, without claiming observational identification of causal effects.

Rather than conceptualizing dropout as a static outcome, we model withdrawal as a dynamic process unfolding over academic time. We operationalize dropout as administrative withdrawal from a course run with an observed unregistration time, enabling a time-to-event representation in discrete academic weeks. The predictive backbone is a discrete-time hazard framework fit on weekly person--period observations, with calibration to yield interpretable hazard probabilities over time.

Beyond prediction, the framework incorporates a \emph{structural counterfactual policy simulation} layer. Under explicit user-defined parameters, model-implied hazards are scaled to represent a hypothetical support regime, enabling comparison of survival trajectories across structural scenarios. In this setting, predicted hazards are interpreted as ``warning windows'' that support sensitivity analysis of intervention timing and magnitude along each enrollment trajectory, while maintaining a clear distinction between regime contrasts in model-implied risk and real-world intervention effects.

The framework \emph{can be} applied in a subgroup-sensitive way for groups defined by observable characteristics. To evaluate this capability, we apply the same simulated regime to observable groups (e.g., disability status) and examine how the regime contrast translates into changes in the survival gap between groups. This analysis does not claim causal identification of subgroup effects; instead, it demonstrates that the mathematical framework supports consistent structural policy evaluation across observable groups, which may motivate future extensions toward more targeted support design.

This study uses the Open University Learning Analytics Dataset (OULAD), which combines time-stamped virtual learning environment (VLE) interaction logs, assessment information, and administrative records \cite{Kuzilek2017OULAD}. The framework is designed to be transferable \emph{when} (i) time-stamped engagement traces and (ii) a well-defined time-to-event endpoint are available, as in many LMS and administrative contexts. In addition, the pipeline includes auxiliary checks---such as calibration/robustness diagnostics and comparison against an alternative baseline---to support empirical validation of the framework.

The paper evaluates the framework through three research questions and their corresponding empirical outputs: (i) construction of the weekly person--period dataset and the temporal split; (ii) temporal risk modeling (hazard) and calibration, producing the baseline risk and implied survival trajectory for \textbf{RQ1}; (iii) structural counterfactual policy simulation (\textbf{RQ2}); and (iv) subgroup validation via gap analysis (\textbf{RQ3}). The next section provides background on student dropout and motivates a temporal, log-based perspective in online learning settings, which sets the stage for the formal definitions and modeling introduced later.